% Generated draft: requires mathematical and source review.
\documentclass[11pt]{article}
\usepackage[margin=25mm]{geometry}
\usepackage{fontspec}
\setmainfont{Latin Modern Roman}
\usepackage{amsmath,amssymb,mathtools,unicode-math}
\setmathfont{Latin Modern Math}
\usepackage{enumitem,needspace,xurl,hyperref}
\setlength{\parindent}{0pt}
\setlength{\parskip}{5pt}
\setlength{\emergencystretch}{3em}
\newcommand{\N}{\mathbb{N}}
\newcommand{\Z}{\mathbb{Z}}
\newcommand{\Q}{\mathbb{Q}}
\newcommand{\R}{\mathbb{R}}
\newcommand{\C}{\mathbb{C}}
\newcommand{\F}{\mathbb{F}}
\newcommand{\A}{\mathbb{A}}
\newcommand{\PP}{\mathbb{P}}
\newcommand{\E}{\mathbb{E}}
\newcommand{\eps}{\varepsilon}
\newcommand{\dd}{\,\mathrm{d}}
\begin{document}
\section*{UnsolvedMath ID: 10400131}
\textbf{Source problem code:} catalogue entry
\subsection{Definitions and mathematical statement}
The problem attributed to S. Morita [228] is to construct a Chern--Simons invariant \(\operatorname{CS}(M)\) as an invariant of the topology of \(M\) itself, for both of the following classes:
\begin{enumerate}
\item every closed oriented three-manifold \(M\);
\item every knot or link complement \(M\) in a closed three-manifold.
\end{enumerate}
It has two expressly stated requirements: give a topological or combinatorial definition, and extend the intended invariant to nonhyperbolic three-manifolds. Thus, after identifying its target \(\mathcal T\) and equivalence convention, the requested assignment must satisfy
\[M\cong M'\ \Longrightarrow\ \operatorname{CS}(M)=\operatorname{CS}(M')\in\mathcal T,\]
with orientation-preserving equivalence in the oriented category, and must recover the source's intended Chern--Simons invariant wherever that is already defined, in particular in its hyperbolic setting. The target, modulus, normalization and the meaning of recovery are not specified here; \(\mathcal T\) is only a placeholder for this missing choice.

For a link \(L\) in a closed manifold \(Y\), ``complement'' could mean \(Y\setminus L\) or the compact exterior \(Y\setminus\operatorname{int}\nu(L)\), where \(\nu(L)\) is a tubular neighborhood. The excerpt does not state which description or which peripheral/boundary data are required. Its second clause does not expressly repeat the orientation assumption, so that must also be clarified. A Chern--Simons functional attached to a chosen connection or extra geometric data alone does not meet the stated goal unless the choice is eliminated or its topological role is specified. No arbitrary connection, metric or representation is selected as a definition in this draft.
\subsection{Short English statement}
Define the intended manifold Chern--Simons invariant topologically or combinatorially for all closed oriented three-manifolds and link complements, including nonhyperbolic cases. The source leaves its normalization and complement boundary conventions unspecified.
\subsection{Sources}
\begin{itemize}
\item Local UnsolvedMath registry, permanent ID \(10400131\), exact source code \texttt{catalogue entry}.
\item S. Morita [228], attribution in Ohtsuki, \emph{Problems on invariants of knots and 3-manifolds} (2002), Problem 7.16, PDF page 111: \url{https://msp.org/gtm/2002/04/gtm-2002-04-024s.pdf}. Full reference [228] is not supplied. Recorded access 2026-07-29; pdf-font-structure extraction.
\end{itemize}
\end{document}
